\documentclass[11pt,a4paper]{article}
\usepackage[utf8]{inputenc}
\usepackage[T1]{fontenc}
\usepackage[margin=0.75in]{geometry}
\usepackage{enumitem}
\usepackage{titlesec}
\usepackage{xcolor}
\usepackage{hyperref}
\definecolor{primarycolor}{RGB}{33,37,41}
\definecolor{secondarycolor}{RGB}{0,102,204}
\definecolor{accentcolor}{RGB}{108,117,125}
\hypersetup{colorlinks=true,urlcolor=secondarycolor}
\titleformat{\section}{\large\bfseries\color{primarycolor}}{}{0em}{}
\titlespacing*{\section}{0pt}{1.2em}{0.5em}
\setlist[itemize]{leftmargin=*,itemsep=0.15em,topsep=0.2em}
\pagestyle{empty}
\begin{document}
\begin{center}
{\Huge\bfseries\color{primarycolor} MARTA KOWAL-BRANDT}\\[0.35em]
{\large QA Automation Engineer | CI/CD \& Platform Engineering}\\[0.25em]
Wroclaw, Poland | marta.kowalbrandt@example.com | +48 555 303 303 | GitLab: gitlab.com/mkb-example
\end{center}
\section*{PROFESSIONAL SUMMARY}
QA automation engineer with experience spanning quality engineering, CI/CD pipelines, test environments and build tooling. Owns GitLab CI pipelines and has delivered faster builds, ephemeral Kubernetes test environments and test observability for multiple product teams. Also builds and monitors a personal k3s homelab, with hands-on Terraform experience provisioning a small Hetzner VPS.
\section*{CORE SKILLS}
CI/CD: GitLab CI, Jenkins, Docker, Docker Compose, Helm\\
Containers and orchestration: Kubernetes (namespace-level access; Helm deployments to shared cluster), k3s (personal homelab)\\
Infrastructure as Code: Terraform (personal Hetzner VPS); Ansible (basic)\\
Cloud: AWS (learning: EC2 and S3 basics)\\
Monitoring and data: Grafana, Prometheus, Loki, PostgreSQL\\
Programming and scripting: Python, Bash, Java (test code), SQL\\
Testing: pytest, Playwright, Selenium, Postman\\
Systems and tools: Linux, Git, Argo CD, Traefik, cert-manager
\section*{PROFESSIONAL EXPERIENCE}
\textbf{Homelab k3s Cluster with GitOps | Personal | 2024-2025}
\begin{itemize}
\item Run a three-node Raspberry Pi 5 k3s cluster, deploying applications from a Git repository with Argo CD.
\item Set up Prometheus, Grafana and Loki monitoring, Traefik ingress and cert-manager for TLS.
\item Provisioned infrastructure for a small Hetzner VPS with Terraform; used Ansible at a basic level.
\end{itemize}
\textbf{QA Automation Engineer (Test Infrastructure) | Odrava Health Software | Wroclaw | 2021-Present}
\begin{itemize}
\item Own GitLab CI pipelines for three product teams (approximately 25 developers), covering builds, unit, integration and end-to-end tests, and artifact publishing.
\item Reduced average pipeline time from 38 to 14 minutes by parallelising test jobs, caching dependencies and separating end-to-end tests into their own stage; monthly CI runner costs fell by about 30\%.
\item Built short-lived merge-request test environments with Docker Compose and later Helm on the shared Kubernetes cluster; environments include the application, PostgreSQL and seeded data, and are automatically removed after merge or 48 hours of inactivity.
\item Maintain a pytest and Playwright end-to-end suite of approximately 900 tests; reduced the flaky-test rate from 7\% to below 1\%.
\item Built Grafana dashboards for flaky tests using CI results, PostgreSQL and the GitLab API.
\item Wrote Bash and Python tools for PostgreSQL test-data seeding and an internal CI guide used in onboarding.
\item Mentored two QA interns in 2023.
\end{itemize}
\textbf{QA Engineer | Brightmill Software House | Wroclaw | 2019-2021}
\begin{itemize}
\item Tested web and mobile applications for client projects using manual and automated approaches.
\item Introduced a Selenium and Java automated regression suite for the largest client project.
\item Configured Jenkins jobs for nightly regression runs.
\end{itemize}
\textbf{Customer Support Specialist (2nd line) | Brightmill Software House | 2018-2019}
\begin{itemize}
\item Investigated customer issues in logs and databases and reproduced bugs for developers.
\end{itemize}
\section*{SELECTED PROJECTS}
\textbf{Merge-Request Ephemeral Environments | Odrava Health | 2023-2024}
\begin{itemize}
\item Deployed a short-lived application and PostgreSQL environment with seeded data for each merge request, using GitLab CI and Helm on the shared Kubernetes cluster.
\item Automatically removed environments after merge or 48 hours of inactivity; enabled QA and product reviews before merge and end-to-end testing against each environment.
\end{itemize}
\textbf{CI Pipeline Speed-up | Odrava Health | 2022}
\begin{itemize}
\item Profiled pipelines for three teams, parallelised test jobs, added dependency caching and separated end-to-end tests into their own stage.
\item Reduced average pipeline time from 38 to 14 minutes and monthly CI runner costs by about 30\%.
\end{itemize}
\textbf{Flaky Test Dashboard | Odrava Health | 2023}
\begin{itemize}
\item Collected CI test results in PostgreSQL and built Grafana dashboards showing flaky tests by team.
\item The flaky-test rate fell from 7\% to below 1\% after teams began addressing the leading causes.
\end{itemize}
\section*{EDUCATION}
BEng, Computer Engineering | Silesian Lowlands University of Technology | 2018
\section*{CERTIFICATIONS AND PROFESSIONAL DEVELOPMENT}
ISTQB Foundation Level | 2019\\
Certified Kubernetes Administrator (CKA) preparation | Exam booked for January 2027; not yet certified
\section*{LANGUAGES}
Polish (native), English (C1), German (A2)
\end{document}